\section{高宗时期的家庭、劳动与物质证据}
高宗时代的国家能力经由赋役、河工、矿业、盐政、驻防、屯田和书籍征集进入不同区域。本节不把户部、方略或制度条文误读为家庭生活的完整档案；它以账本所列材料为线索，说明何处可以讨论劳动、迁徙、教育、宗教与物质供给，何处只能承认材料不足。
\subsection{1720—1729：地方材料所能看见的尺度}
以下事件按账本时间排列；它们不是同质的社会样本，而是提示应到何种地点、机构和材料中继续追问的入口。
\hypertarget{social:EQ-1723-YONGZHENG-ACCESSION}{}\paragraph{雍正元年（1723年）}雍正帝即位并开始重整康熙晚年的财政、宗室与官僚议题。从地方治理、身份、家庭与物质生活的角度，这一节点要求追问它在何处改变了资源、道路、组织或日常风险。核心来源为《清世宗实录》雍正元年相关条；宫中朱批奏折、内阁大库题本与《清史稿》相关志传。它首先留下的是官府分类、任命、处置或资源调拨的痕迹；要判断基层是否照办，仍须寻找县档、账簿、契约或后续案件。账本所示的研究方向是宫廷政治、制度运作与中央—地方信息链研究。可核的条件包括康熙帝去世后继承程序与皇子网络需要迅速稳定。应分层观察的后果是新帝以密折、用人与财政整顿重塑决策中心。证据界限：以雍正元年（1723年）的同期文书为定位起点；官修记录、奏报或外交文本服务特定机构立场，具体日期、规模、地方执行与对手视角须以独立材料复核。
\hypertarget{social:EQ-1723-JI-ZENGYUN-CASE}{}\paragraph{雍正元年（1723年）}年羹尧在西北军政中的权力与朝廷猜忌加剧。从地方治理、身份、家庭与物质生活的角度，这一节点要求追问它在何处改变了资源、道路、组织或日常风险。核心来源为《清世宗实录》雍正元年相关条；宫中朱批奏折、内阁大库题本与《清史稿》相关志传。它首先留下的是官府分类、任命、处置或资源调拨的痕迹；要判断基层是否照办，仍须寻找县档、账簿、契约或后续案件。账本所示的研究方向是宫廷政治、制度运作与中央—地方信息链研究。可核的条件包括边军统帅兼具军功、财政与人事资源。应分层观察的后果是其后对年氏集团的处分成为雍正整顿权力的节点。证据界限：以雍正元年（1723年）的同期文书为定位起点；官修记录、奏报或外交文本服务特定机构立场，具体日期、规模、地方执行与对手视角须以独立材料复核。
\hypertarget{social:EQ-1724-CATHOLIC-PROSCRIPTION}{}\paragraph{雍正二年（1724年）}清廷扩大对天主教活动的限制，将传教、地方教民与外来人员置入治理问题。从地方治理、身份、家庭与物质生活的角度，这一节点要求追问它在何处改变了资源、道路、组织或日常风险。核心来源为《清世宗实录》雍正二年相关条；地方志、督抚奏折及地方档案。这类编年材料可以标定朝廷获知和叙述事件的时间，却往往压低妇女、儿童、佃户、流民和非文书精英的行动。账本所示的研究方向是地方档案、迁徙、宗教、族群与社会动员研究。可核的条件包括礼仪争论与地方官对宗教网络的疑虑相互叠加。应分层观察的后果是宗教实践转入不同地区和身份条件下的隐蔽或协商状态。证据界限：以雍正二年（1724年）的同期文书为定位起点；官修记录、奏报或外交文本服务特定机构立场，具体日期、规模、地方执行与对手视角须以独立材料复核。
\hypertarget{social:EQ-1725-NIAN-GENGYAO-PUNISHMENT}{}\paragraph{雍正三年（1725年）}年羹尧被削职、赐死，其党羽与相关官员受到清理。从地方治理、身份、家庭与物质生活的角度，这一节点要求追问它在何处改变了资源、道路、组织或日常风险。核心来源为《清世宗实录》雍正三年相关条；宫中朱批奏折、内阁大库题本与《清史稿》相关志传。它首先留下的是官府分类、任命、处置或资源调拨的痕迹；要判断基层是否照办，仍须寻找县档、账簿、契约或后续案件。账本所示的研究方向是宫廷政治、制度运作与中央—地方信息链研究。可核的条件包括雍正帝须限制功臣和边帅形成独立政治中心。应分层观察的后果是中央对军政人事和密折信息链的控制增强。证据界限：以雍正三年（1725年）的同期文书为定位起点；官修记录、奏报或外交文本服务特定机构立场，具体日期、规模、地方执行与对手视角须以独立材料复核。
\hypertarget{social:EQ-1725-TAN-TEBI-CASE}{}\paragraph{雍正三年（1725年）}隆科多失势并被幽禁，雍正初年近臣权力结构继续调整。从地方治理、身份、家庭与物质生活的角度，这一节点要求追问它在何处改变了资源、道路、组织或日常风险。核心来源为《清世宗实录》雍正三年相关条；宫中朱批奏折、内阁大库题本与《清史稿》相关志传。它首先留下的是官府分类、任命、处置或资源调拨的痕迹；要判断基层是否照办，仍须寻找县档、账簿、契约或后续案件。账本所示的研究方向是宫廷政治、制度运作与中央—地方信息链研究。可核的条件包括继承过程中的近臣角色引发君主对安全与忠诚的担忧。应分层观察的后果是内廷、宗室与部院之间的政治边界被重新划定。证据界限：以雍正三年（1725年）的同期文书为定位起点；官修记录、奏报或外交文本服务特定机构立场，具体日期、规模、地方执行与对手视角须以独立材料复核。
\hypertarget{social:EQ-1726-HAOXIAN-GUIGONG}{}\paragraph{雍正四年（1726年）}雍正政府推进耗羡归公与养廉银安排，试图重组地方财政和官员俸给。从地方治理、身份、家庭与物质生活的角度，这一节点要求追问它在何处改变了资源、道路、组织或日常风险。核心来源为《清世宗实录》雍正四年相关条；户部奏销、盐政、河工或海关档案。它首先留下的是官府分类、任命、处置或资源调拨的痕迹；要判断基层是否照办，仍须寻找县档、账簿、契约或后续案件。账本所示的研究方向是财政账册、市场网络、灾害环境与区域经济研究。可核的条件包括正额俸银不足与附加征收长期并存。应分层观察的后果是地方财政的公开名目扩大，但实际征收和地区差异仍待账册核验。证据界限：以雍正四年（1726年）的同期文书为定位起点；官修记录、奏报或外交文本服务特定机构立场，具体日期、规模、地方执行与对手视角须以独立材料复核。
\hypertarget{social:EQ-1726-Gaitu-Guiliu-SOUTHWEST}{}\paragraph{雍正四年（1726年）}清廷在西南继续改土归流，将部分土司辖地改置流官体系。从地方治理、身份、家庭与物质生活的角度，这一节点要求追问它在何处改变了资源、道路、组织或日常风险。核心来源为《清世宗实录》雍正四年相关条；军机处档、理藩院档或相关方略。这类编年材料可以标定朝廷获知和叙述事件的时间，却往往压低妇女、儿童、佃户、流民和非文书精英的行动。账本所示的研究方向是多语边疆档案、盟旗/寺院/绿洲地方史与帝国治理研究。可核的条件包括矿产、交通、纠纷与军事控制推动中央介入。应分层观察的后果是行政改置伴随地方抵抗、迁徙和长期治理成本。证据界限：以雍正四年（1726年）的同期文书为定位起点；官修记录、奏报或外交文本服务特定机构立场，具体日期、规模、地方执行与对手视角须以独立材料复核。
\hypertarget{social:EQ-1727-KYAKHTA-TREATY}{}\paragraph{雍正五年（1727年）}《恰克图条约》确定清俄北部边界与贸易、使团安排。从地方治理、身份、家庭与物质生活的角度，这一节点要求追问它在何处改变了资源、道路、组织或日常风险。核心来源为《清世宗实录》雍正五年相关条；《筹办夷务始末》相关年卷与中外照会、条约文本。它首先留下的是官府分类、任命、处置或资源调拨的痕迹；要判断基层是否照办，仍须寻找县档、账簿、契约或后续案件。账本所示的研究方向是条约文本、外交档案、口岸社会与国际法实践研究。可核的条件包括边界冲突和蒙古地区秩序需要制度化协商。应分层观察的后果是恰克图成为跨境贸易与外交往来的重要节点。证据界限：以雍正五年（1727年）的同期文书为定位起点；官修记录、奏报或外交文本服务特定机构立场，具体日期、规模、地方执行与对手视角须以独立材料复核。
\hypertarget{social:EQ-1727-BANNER-GENEALOGIES}{}\paragraph{雍正五年（1727年）}清廷命令编修八旗谱牒，试图界定旗籍、世系与待遇资格。从地方治理、身份、家庭与物质生活的角度，这一节点要求追问它在何处改变了资源、道路、组织或日常风险。核心来源为《清世宗实录》雍正五年相关条；地方志、督抚奏折及地方档案。这类编年材料可以标定朝廷获知和叙述事件的时间，却往往压低妇女、儿童、佃户、流民和非文书精英的行动。账本所示的研究方向是地方档案、迁徙、宗教、族群与社会动员研究。可核的条件包括旗籍资源、婚姻与身份管理需要可查的分类文书。应分层观察的后果是旗人身份被更紧密地纳入国家档案和待遇秩序。证据界限：以雍正五年（1727年）的同期文书为定位起点；官修记录、奏报或外交文本服务特定机构立场，具体日期、规模、地方执行与对手视角须以独立材料复核。
\hypertarget{social:EQ-1728-Dzungar-EXPEDITION}{}\paragraph{雍正六年（1728年）}清军在西北对准噶尔展开远征，后勤与高原草原路线成为军事难题。从地方治理、身份、家庭与物质生活的角度，这一节点要求追问它在何处改变了资源、道路、组织或日常风险。核心来源为《清世宗实录》雍正六年相关条；军机处档、战事方略及地方奏报。这类编年材料可以标定朝廷获知和叙述事件的时间，却往往压低妇女、儿童、佃户、流民和非文书精英的行动。账本所示的研究方向是战争动员、军需、战场暴力与地方社会研究。可核的条件包括清廷希望限制准噶尔对蒙古与青海的影响。应分层观察的后果是军需困难促使中央反复调整将领、补给和联盟策略。证据界限：以雍正六年（1728年）的同期文书为定位起点；官修记录、奏报或外交文本服务特定机构立场，具体日期、规模、地方执行与对手视角须以独立材料复核。
\hypertarget{social:EQ-1728-ZHUNGHAR-TIBET-SETTLEMENT}{}\paragraph{雍正六年（1728年）}清廷在准噶尔威胁下调整驻藏与青海蒙古的军事部署。从地方治理、身份、家庭与物质生活的角度，这一节点要求追问它在何处改变了资源、道路、组织或日常风险。核心来源为《清世宗实录》雍正六年相关条；军机处档、理藩院档或相关方略。这类编年材料可以标定朝廷获知和叙述事件的时间，却往往压低妇女、儿童、佃户、流民和非文书精英的行动。账本所示的研究方向是多语边疆档案、盟旗/寺院/绿洲地方史与帝国治理研究。可核的条件包括拉萨政局和草原力量关系持续不稳。应分层观察的后果是西藏事务被更直接地纳入清廷边务协调。证据界限：以雍正六年（1728年）的同期文书为定位起点；官修记录、奏报或外交文本服务特定机构立场，具体日期、规模、地方执行与对手视角须以独立材料复核。
\hypertarget{social:EQ-1729-JUNJICHU-FOUNDING}{}\paragraph{雍正七年（1729年）}雍正帝设军机房，后来发展为军机处，处理紧急军政文书。从地方治理、身份、家庭与物质生活的角度，这一节点要求追问它在何处改变了资源、道路、组织或日常风险。核心来源为《清世宗实录》雍正七年相关条；宫中朱批奏折、内阁大库题本与《清史稿》相关志传。它首先留下的是官府分类、任命、处置或资源调拨的痕迹；要判断基层是否照办，仍须寻找县档、账簿、契约或后续案件。账本所示的研究方向是宫廷政治、制度运作与中央—地方信息链研究。可核的条件包括西北战争要求比常规部院更快的保密决策链。应分层观察的后果是军机处成为中后期中央决策的重要枢纽。证据界限：以雍正七年（1729年）的同期文书为定位起点；官修记录、奏报或外交文本服务特定机构立场，具体日期、规模、地方执行与对手视角须以独立材料复核。
\subsection{1730—1739：地方材料所能看见的尺度}
以下事件按账本时间排列；它们不是同质的社会样本，而是提示应到何种地点、机构和材料中继续追问的入口。
\hypertarget{social:EQ-1730-SOUTHWEST-MIAO-GOVERNANCE}{}\paragraph{雍正八年（1730年）}清廷围绕苗疆设置、屯兵与地方纠纷加强治理。从地方治理、身份、家庭与物质生活的角度，这一节点要求追问它在何处改变了资源、道路、组织或日常风险。核心来源为《清世宗实录》雍正八年相关条；军机处档、理藩院档或相关方略。这类编年材料可以标定朝廷获知和叙述事件的时间，却往往压低妇女、儿童、佃户、流民和非文书精英的行动。账本所示的研究方向是多语边疆档案、盟旗/寺院/绿洲地方史与帝国治理研究。可核的条件包括改土归流与移民开发改变既有地方权力关系。应分层观察的后果是军政介入与社区抵抗交织，不能以行政设置等同稳定统治。证据界限：以雍正八年（1730年）的同期文书为定位起点；官修记录、奏报或外交文本服务特定机构立场，具体日期、规模、地方执行与对手视角须以独立材料复核。
\hypertarget{social:EQ-1731-Dzungar-WAR-REVERSE}{}\paragraph{雍正九年（1731年）}清军在对准噶尔战争中受挫，西北远征暴露补给与指挥困难。从地方治理、身份、家庭与物质生活的角度，这一节点要求追问它在何处改变了资源、道路、组织或日常风险。核心来源为《清世宗实录》雍正九年相关条；军机处档、战事方略及地方奏报。这类编年材料可以标定朝廷获知和叙述事件的时间，却往往压低妇女、儿童、佃户、流民和非文书精英的行动。账本所示的研究方向是战争动员、军需、战场暴力与地方社会研究。可核的条件包括长距离运输、气候和草原盟友关系限制军队行动。应分层观察的后果是朝廷调整统帅、粮道与军费安排。证据界限：以雍正九年（1731年）的同期文书为定位起点；官修记录、奏报或外交文本服务特定机构立场，具体日期、规模、地方执行与对手视角须以独立材料复核。
\hypertarget{social:EQ-1732-KOBDO-BATTLE}{}\paragraph{雍正十年（1732年）}清军在和通泊等战事失利，西北战局一度紧张。从地方治理、身份、家庭与物质生活的角度，这一节点要求追问它在何处改变了资源、道路、组织或日常风险。核心来源为《清世宗实录》雍正十年相关条；军机处档、战事方略及地方奏报。这类编年材料可以标定朝廷获知和叙述事件的时间，却往往压低妇女、儿童、佃户、流民和非文书精英的行动。账本所示的研究方向是战争动员、军需、战场暴力与地方社会研究。可核的条件包括准噶尔骑兵和草原地理使清军难以速胜。应分层观察的后果是清廷转向重整蒙古盟友、驻防与运输体系。证据界限：以雍正十年（1732年）的同期文书为定位起点；官修记录、奏报或外交文本服务特定机构立场，具体日期、规模、地方执行与对手视角须以独立材料复核。
\hypertarget{social:EQ-1732-AMIN-KHOJA-FLIGHT}{}\paragraph{雍正十年（1732年）}阿敏和卓等人在准噶尔压力下向清境迁移，绿洲政治进入清准竞争。从地方治理、身份、家庭与物质生活的角度，这一节点要求追问它在何处改变了资源、道路、组织或日常风险。核心来源为《清世宗实录》雍正十年相关条；军机处档、理藩院档或相关方略。这类编年材料可以标定朝廷获知和叙述事件的时间，却往往压低妇女、儿童、佃户、流民和非文书精英的行动。账本所示的研究方向是多语边疆档案、盟旗/寺院/绿洲地方史与帝国治理研究。可核的条件包括准噶尔内部统治与绿洲地方关系紧张。应分层观察的后果是清廷获得介入天山南路事务的新信息和人际网络。证据界限：以雍正十年（1732年）的同期文书为定位起点；官修记录、奏报或外交文本服务特定机构立场，具体日期、规模、地方执行与对手视角须以独立材料复核。
\hypertarget{social:EQ-1733-YONGZHENG-FISCAL-INSPECTION}{}\paragraph{雍正十一年（1733年）}雍正政府持续以奏销、亏空追查和官员考成为手段整顿地方财政。从地方治理、身份、家庭与物质生活的角度，这一节点要求追问它在何处改变了资源、道路、组织或日常风险。核心来源为《清世宗实录》雍正十一年相关条；户部奏销、盐政、河工或海关档案。它首先留下的是官府分类、任命、处置或资源调拨的痕迹；要判断基层是否照办，仍须寻找县档、账簿、契约或后续案件。账本所示的研究方向是财政账册、市场网络、灾害环境与区域经济研究。可核的条件包括耗羡改革需要与旧有钱粮、仓储和欠额制度衔接。应分层观察的后果是地方官的呈报策略与财政实际之间形成新的张力。证据界限：以雍正十一年（1733年）的同期文书为定位起点；官修记录、奏报或外交文本服务特定机构立场，具体日期、规模、地方执行与对手视角须以独立材料复核。
\hypertarget{social:EQ-1734-CHRISTIAN-ENFORCEMENT}{}\paragraph{雍正十二年（1734年）}地方官继续执行限制天主教的政策，传教士和教民的活动受地域差异影响。从地方治理、身份、家庭与物质生活的角度，这一节点要求追问它在何处改变了资源、道路、组织或日常风险。核心来源为《清世宗实录》雍正十二年相关条；地方志、督抚奏折及地方档案。这类编年材料可以标定朝廷获知和叙述事件的时间，却往往压低妇女、儿童、佃户、流民和非文书精英的行动。账本所示的研究方向是地方档案、迁徙、宗教、族群与社会动员研究。可核的条件包括中央禁令需经省县转化为具体案件与登记。应分层观察的后果是不同地区的执行强度不能由中央谕旨一概而论。证据界限：以雍正十二年（1734年）的同期文书为定位起点；官修记录、奏报或外交文本服务特定机构立场，具体日期、规模、地方执行与对手视角须以独立材料复核。
\hypertarget{social:EQ-1735-QIANLONG-ACCESSION}{}\paragraph{雍正十三年（1735年）}乾隆帝即位，雍正时期形成的财政与军机决策机制进入新朝。从地方治理、身份、家庭与物质生活的角度，这一节点要求追问它在何处改变了资源、道路、组织或日常风险。核心来源为《清世宗实录》雍正十三年相关条；宫中朱批奏折、内阁大库题本与《清史稿》相关志传。它首先留下的是官府分类、任命、处置或资源调拨的痕迹；要判断基层是否照办，仍须寻找县档、账簿、契约或后续案件。账本所示的研究方向是宫廷政治、制度运作与中央—地方信息链研究。可核的条件包括雍正帝去世与皇位继承完成。应分层观察的后果是乾隆初政开始重估前朝用人、刑狱与边务方针。证据界限：以雍正十三年（1735年）的同期文书为定位起点；官修记录、奏报或外交文本服务特定机构立场，具体日期、规模、地方执行与对手视角须以独立材料复核。
\hypertarget{social:EQ-1735-MIAO-UPRISING-RONGJIANG}{}\paragraph{雍正十三年（1735年）}贵州荣江等地苗侗社区反抗清军和地方治理，镇压造成严重暴力。从地方治理、身份、家庭与物质生活的角度，这一节点要求追问它在何处改变了资源、道路、组织或日常风险。核心来源为《清世宗实录》雍正十三年相关条；军机处档、战事方略及地方奏报。这类编年材料可以标定朝廷获知和叙述事件的时间，却往往压低妇女、儿童、佃户、流民和非文书精英的行动。账本所示的研究方向是战争动员、军需、战场暴力与地方社会研究。可核的条件包括行政扩张、赋役与地方资源冲突累积。应分层观察的后果是官方战报的伤亡数字和族群标签须与地方材料互校。证据界限：以雍正十三年（1735年）的同期文书为定位起点；官修记录、奏报或外交文本服务特定机构立场，具体日期、规模、地方执行与对手视角须以独立材料复核。
\hypertarget{social:EQ-1736-AMNESTY-AND-RECTIFICATION}{}\paragraph{乾隆元年（1736年）}乾隆初年颁布恩诏并调整雍正后期政治案件的处置。从地方治理、身份、家庭与物质生活的角度，这一节点要求追问它在何处改变了资源、道路、组织或日常风险。核心来源为《清高宗实录》乾隆元年相关条；宫中朱批奏折、内阁大库题本与《清史稿》相关志传。它首先留下的是官府分类、任命、处置或资源调拨的痕迹；要判断基层是否照办，仍须寻找县档、账簿、契约或后续案件。账本所示的研究方向是宫廷政治、制度运作与中央—地方信息链研究。可核的条件包括新君需要建立自身施政形象并重组官僚关系。应分层观察的后果是部分案件被重新解释，但中央集权并未放松。证据界限：以乾隆元年（1736年）的同期文书为定位起点；官修记录、奏报或外交文本服务特定机构立场，具体日期、规模、地方执行与对手视角须以独立材料复核。
\hypertarget{social:EQ-1737-Dzungar-NEGOTIATION}{}\paragraph{乾隆二年（1737年）}清廷在战争压力下继续处理与准噶尔的使节、盟旗和军事情报。从地方治理、身份、家庭与物质生活的角度，这一节点要求追问它在何处改变了资源、道路、组织或日常风险。核心来源为《清高宗实录》乾隆二年相关条；《筹办夷务始末》相关年卷与中外照会、条约文本。它首先留下的是官府分类、任命、处置或资源调拨的痕迹；要判断基层是否照办，仍须寻找县档、账簿、契约或后续案件。账本所示的研究方向是条约文本、外交档案、口岸社会与国际法实践研究。可核的条件包括草原战争难以只靠单线军事行动解决。应分层观察的后果是外交、情报与军事准备并行发展。证据界限：以乾隆二年（1737年）的同期文书为定位起点；官修记录、奏报或外交文本服务特定机构立场，具体日期、规模、地方执行与对手视角须以独立材料复核。
\hypertarget{social:EQ-1738-FIRST-JINCHUAN-WAR}{}\paragraph{乾隆三年（1738年）}清军开始大规模干预大金川，山地战与土司关系限制军事推进。从地方治理、身份、家庭与物质生活的角度，这一节点要求追问它在何处改变了资源、道路、组织或日常风险。核心来源为《清高宗实录》乾隆三年相关条；军机处档、战事方略及地方奏报。这类编年材料可以标定朝廷获知和叙述事件的时间，却往往压低妇女、儿童、佃户、流民和非文书精英的行动。账本所示的研究方向是战争动员、军需、战场暴力与地方社会研究。可核的条件包括地方权力冲突与清廷对川西治理的扩展相互作用。应分层观察的后果是远征显示高山补给和地方联盟的成本极高。证据界限：以乾隆三年（1738年）的同期文书为定位起点；官修记录、奏报或外交文本服务特定机构立场，具体日期、规模、地方执行与对手视角须以独立材料复核。
\hypertarget{social:EQ-1739-JINCHUAN-SUPPLY}{}\paragraph{乾隆四年（1739年）}大金川战事的粮运、道路和军费成为中央持续关注的难题。从地方治理、身份、家庭与物质生活的角度，这一节点要求追问它在何处改变了资源、道路、组织或日常风险。核心来源为《清高宗实录》乾隆四年相关条；户部奏销、盐政、河工或海关档案。它首先留下的是官府分类、任命、处置或资源调拨的痕迹；要判断基层是否照办，仍须寻找县档、账簿、契约或后续案件。账本所示的研究方向是财政账册、市场网络、灾害环境与区域经济研究。可核的条件包括山地驻军需要远距离转运粮械。应分层观察的后果是军费和民夫负担使战争后果超出战场范围。证据界限：以乾隆四年（1739年）的同期文书为定位起点；官修记录、奏报或外交文本服务特定机构立场，具体日期、规模、地方执行与对手视角须以独立材料复核。
\subsection{1740—1749：地方材料所能看见的尺度}
以下事件按账本时间排列；它们不是同质的社会样本，而是提示应到何种地点、机构和材料中继续追问的入口。
\hypertarget{social:EQ-1740-NORTH-CHINA-FAMINE}{}\paragraph{乾隆五年（1740年）}华北旱灾和粮价问题促使地方奏报与赈济措施增加。从地方治理、身份、家庭与物质生活的角度，这一节点要求追问它在何处改变了资源、道路、组织或日常风险。核心来源为《清高宗实录》乾隆五年相关条；地方志、督抚奏折及地方档案。这类编年材料可以标定朝廷获知和叙述事件的时间，却往往压低妇女、儿童、佃户、流民和非文书精英的行动。账本所示的研究方向是地方档案、迁徙、宗教、族群与社会动员研究。可核的条件包括气候波动、粮食流通和地方仓储共同影响灾情。应分层观察的后果是灾情、赈济和死亡迁徙应分开记录并核对地域口径。证据界限：以乾隆五年（1740年）的同期文书为定位起点；官修记录、奏报或外交文本服务特定机构立场，具体日期、规模、地方执行与对手视角须以独立材料复核。
\hypertarget{social:EQ-1741-MIAO-FRONTIER-CAMPAIGN}{}\paragraph{乾隆六年（1741年）}清廷在西南继续以驻兵、设治和军事行动处理苗疆冲突。从地方治理、身份、家庭与物质生活的角度，这一节点要求追问它在何处改变了资源、道路、组织或日常风险。核心来源为《清高宗实录》乾隆六年相关条；军机处档、理藩院档或相关方略。这类编年材料可以标定朝廷获知和叙述事件的时间，却往往压低妇女、儿童、佃户、流民和非文书精英的行动。账本所示的研究方向是多语边疆档案、盟旗/寺院/绿洲地方史与帝国治理研究。可核的条件包括地方行政延伸与社区自治之间矛盾未解。应分层观察的后果是长期驻防和地方中介成为治理的重要条件。证据界限：以乾隆六年（1741年）的同期文书为定位起点；官修记录、奏报或外交文本服务特定机构立场，具体日期、规模、地方执行与对手视角须以独立材料复核。
\hypertarget{social:EQ-1742-BANNER-EXIT-ORDER}{}\paragraph{乾隆七年（1742年）}清廷允许部分汉军旗人出旗，调整旗籍与生计安排。从地方治理、身份、家庭与物质生活的角度，这一节点要求追问它在何处改变了资源、道路、组织或日常风险。核心来源为《清高宗实录》乾隆七年相关条；地方志、督抚奏折及地方档案。这类编年材料可以标定朝廷获知和叙述事件的时间，却往往压低妇女、儿童、佃户、流民和非文书精英的行动。账本所示的研究方向是地方档案、迁徙、宗教、族群与社会动员研究。可核的条件包括八旗人口增长和财政供养压力加重。应分层观察的后果是出旗改变身份和待遇资格，实际执行具有地区差异。证据界限：以乾隆七年（1742年）的同期文书为定位起点；官修记录、奏报或外交文本服务特定机构立场，具体日期、规模、地方执行与对手视角须以独立材料复核。
\hypertarget{social:EQ-1743-YELLOW-RIVER-WORKS}{}\paragraph{乾隆八年（1743年）}河工官员就黄河堤防、河道和经费持续奏报并实施工程。从地方治理、身份、家庭与物质生活的角度，这一节点要求追问它在何处改变了资源、道路、组织或日常风险。核心来源为《清高宗实录》乾隆八年相关条；户部奏销、盐政、河工或海关档案。它首先留下的是官府分类、任命、处置或资源调拨的痕迹；要判断基层是否照办，仍须寻找县档、账簿、契约或后续案件。账本所示的研究方向是财政账册、市场网络、灾害环境与区域经济研究。可核的条件包括河患威胁漕运、田赋与沿河社区。应分层观察的后果是工程预算与实际河势、民工负担不能由定额直接推断。证据界限：以乾隆八年（1743年）的同期文书为定位起点；官修记录、奏报或外交文本服务特定机构立场，具体日期、规模、地方执行与对手视角须以独立材料复核。
\hypertarget{social:EQ-1744-JESUIT-CALENDAR-SERVICE}{}\paragraph{乾隆九年（1744年）}宫廷继续利用部分西洋传教士的历算与技术服务，同时维持宗教限制。从地方治理、身份、家庭与物质生活的角度，这一节点要求追问它在何处改变了资源、道路、组织或日常风险。核心来源为《清高宗实录》乾隆九年相关条；内阁档、文集或报刊相关记录。这类编年材料可以标定朝廷获知和叙述事件的时间，却往往压低妇女、儿童、佃户、流民和非文书精英的行动。账本所示的研究方向是知识生产、教育、宗教与公共传播研究。可核的条件包括历法、仪器与宫廷知识需求并未与宗教政策完全重合。应分层观察的后果是技术合作与传教许可须分开考察。证据界限：以乾隆九年（1744年）的同期文书为定位起点；官修记录、奏报或外交文本服务特定机构立场，具体日期、规模、地方执行与对手视角须以独立材料复核。
\hypertarget{social:EQ-1745-JINCHUAN-NEGOTIATION}{}\paragraph{乾隆十年（1745年）}清廷在大金川战局中尝试通过招抚、盟约和军事压力重组地方关系。从地方治理、身份、家庭与物质生活的角度，这一节点要求追问它在何处改变了资源、道路、组织或日常风险。核心来源为《清高宗实录》乾隆十年相关条；军机处档、理藩院档或相关方略。这类编年材料可以标定朝廷获知和叙述事件的时间，却往往压低妇女、儿童、佃户、流民和非文书精英的行动。账本所示的研究方向是多语边疆档案、盟旗/寺院/绿洲地方史与帝国治理研究。可核的条件包括长期作战难以迅速取得山地控制。应分层观察的后果是地方首领的降附文书不等于基层社会立即服从。证据界限：以乾隆十年（1745年）的同期文书为定位起点；官修记录、奏报或外交文本服务特定机构立场，具体日期、规模、地方执行与对手视角须以独立材料复核。
\hypertarget{social:EQ-1746-FIRST-JINCHUAN-SETTLEMENT}{}\paragraph{乾隆十一年（1746年）}第一次金川战争以和议和行政安排告一段落。从地方治理、身份、家庭与物质生活的角度，这一节点要求追问它在何处改变了资源、道路、组织或日常风险。核心来源为《清高宗实录》乾隆十一年相关条；军机处档、战事方略及地方奏报。这类编年材料可以标定朝廷获知和叙述事件的时间，却往往压低妇女、儿童、佃户、流民和非文书精英的行动。账本所示的研究方向是战争动员、军需、战场暴力与地方社会研究。可核的条件包括军费高昂与战场推进困难促使朝廷接受阶段性解决。应分层观察的后果是后续土司关系并未根除，为再战留下条件。证据界限：以乾隆十一年（1746年）的同期文书为定位起点；官修记录、奏报或外交文本服务特定机构立场，具体日期、规模、地方执行与对手视角须以独立材料复核。
\hypertarget{social:EQ-1747-SOUTHWEST-MINING-REGULATION}{}\paragraph{乾隆十二年（1747年）}清廷和地方官围绕银铜矿开采、税课与矿工流动加强管理。从地方治理、身份、家庭与物质生活的角度，这一节点要求追问它在何处改变了资源、道路、组织或日常风险。核心来源为《清高宗实录》乾隆十二年相关条；户部奏销、盐政、河工或海关档案。它首先留下的是官府分类、任命、处置或资源调拨的痕迹；要判断基层是否照办，仍须寻找县档、账簿、契约或后续案件。账本所示的研究方向是财政账册、市场网络、灾害环境与区域经济研究。可核的条件包括矿业吸引流民、资本和武装群体，治安风险上升。应分层观察的后果是矿务文书可见国家规制，不能代表全部地下经济。证据界限：以乾隆十二年（1747年）的同期文书为定位起点；官修记录、奏报或外交文本服务特定机构立场，具体日期、规模、地方执行与对手视角须以独立材料复核。
\hypertarget{social:EQ-1748-SONG-AND-CASE}{}\paragraph{乾隆十三年（1748年）}乾隆朝通过大案处分官员，强化对地方文书、钱粮和吏治的监督。从地方治理、身份、家庭与物质生活的角度，这一节点要求追问它在何处改变了资源、道路、组织或日常风险。核心来源为《清高宗实录》乾隆十三年相关条；宫中朱批奏折、内阁大库题本与《清史稿》相关志传。它首先留下的是官府分类、任命、处置或资源调拨的痕迹；要判断基层是否照办，仍须寻找县档、账簿、契约或后续案件。账本所示的研究方向是宫廷政治、制度运作与中央—地方信息链研究。可核的条件包括中央担忧地方官隐瞒亏空与相互包庇。应分层观察的后果是案件成为考成与惩戒制度的示范，但不证明全面廉洁。证据界限：以乾隆十三年（1748年）的同期文书为定位起点；官修记录、奏报或外交文本服务特定机构立场，具体日期、规模、地方执行与对手视角须以独立材料复核。
\hypertarget{social:EQ-1749-JIANGNAN-TAX-REVIEW}{}\paragraph{乾隆十四年（1749年）}江南钱粮、漕运和士绅关系继续在户部与地方奏折中被讨论。从地方治理、身份、家庭与物质生活的角度，这一节点要求追问它在何处改变了资源、道路、组织或日常风险。核心来源为《清高宗实录》乾隆十四年相关条；户部奏销、盐政、河工或海关档案。它首先留下的是官府分类、任命、处置或资源调拨的痕迹；要判断基层是否照办，仍须寻找县档、账簿、契约或后续案件。账本所示的研究方向是财政账册、市场网络、灾害环境与区域经济研究。可核的条件包括富庶地区的税粮征解和地方负担具有全国财政意义。应分层观察的后果是中央定额与地方实际征收之间的差距需依账册检验。证据界限：以乾隆十四年（1749年）的同期文书为定位起点；官修记录、奏报或外交文本服务特定机构立场，具体日期、规模、地方执行与对手视角须以独立材料复核。
\subsection{1750—1759：地方材料所能看见的尺度}
以下事件按账本时间排列；它们不是同质的社会样本，而是提示应到何种地点、机构和材料中继续追问的入口。
\hypertarget{social:EQ-1750-LHASA-UPRISING}{}\paragraph{乾隆十五年（1750年）}拉萨发生政治冲突，驻藏大臣和清军介入处置。从地方治理、身份、家庭与物质生活的角度，这一节点要求追问它在何处改变了资源、道路、组织或日常风险。核心来源为《清高宗实录》乾隆十五年相关条；军机处档、理藩院档或相关方略。这类编年材料可以标定朝廷获知和叙述事件的时间，却往往压低妇女、儿童、佃户、流民和非文书精英的行动。账本所示的研究方向是多语边疆档案、盟旗/寺院/绿洲地方史与帝国治理研究。可核的条件包括西藏地方权力、驻防与宗教政治关系不稳。应分层观察的后果是清廷随后调整驻藏制度，但不能抹去地方行动者的作用。证据界限：以乾隆十五年（1750年）的同期文书为定位起点；官修记录、奏报或外交文本服务特定机构立场，具体日期、规模、地方执行与对手视角须以独立材料复核。
\hypertarget{social:EQ-1751-TIBET-ORDINANCE}{}\paragraph{乾隆十六年（1751年）}清廷在拉萨事变后重整驻藏大臣和噶厦相关制度安排。从地方治理、身份、家庭与物质生活的角度，这一节点要求追问它在何处改变了资源、道路、组织或日常风险。核心来源为《清高宗实录》乾隆十六年相关条；军机处档、理藩院档或相关方略。这类编年材料可以标定朝廷获知和叙述事件的时间，却往往压低妇女、儿童、佃户、流民和非文书精英的行动。账本所示的研究方向是多语边疆档案、盟旗/寺院/绿洲地方史与帝国治理研究。可核的条件包括1750年冲突暴露既有权力安排的脆弱性。应分层观察的后果是制度条文与西藏地方实际政治运作需要并读藏汉材料。证据界限：以乾隆十六年（1751年）的同期文书为定位起点；官修记录、奏报或外交文本服务特定机构立场，具体日期、规模、地方执行与对手视角须以独立材料复核。
\hypertarget{social:EQ-1752-TIBETAN-ADMINISTRATION}{}\paragraph{乾隆十七年（1752年）}驻藏系统与青海、川藏交通的军政协调继续调整。从地方治理、身份、家庭与物质生活的角度，这一节点要求追问它在何处改变了资源、道路、组织或日常风险。核心来源为《清高宗实录》乾隆十七年相关条；军机处档、理藩院档或相关方略。这类编年材料可以标定朝廷获知和叙述事件的时间，却往往压低妇女、儿童、佃户、流民和非文书精英的行动。账本所示的研究方向是多语边疆档案、盟旗/寺院/绿洲地方史与帝国治理研究。可核的条件包括高原补给和多方政治关系限制中央命令的执行。应分层观察的后果是边务档案可显示决策链，不应替代地方社会证词。证据界限：以乾隆十七年（1752年）的同期文书为定位起点；官修记录、奏报或外交文本服务特定机构立场，具体日期、规模、地方执行与对手视角须以独立材料复核。
\hypertarget{social:EQ-1753-SALT-ADMINISTRATION}{}\paragraph{乾隆十八年（1753年）}盐政、引岸和课额问题在中央与产销地区持续受到整顿。从地方治理、身份、家庭与物质生活的角度，这一节点要求追问它在何处改变了资源、道路、组织或日常风险。核心来源为《清高宗实录》乾隆十八年相关条；户部奏销、盐政、河工或海关档案。它首先留下的是官府分类、任命、处置或资源调拨的痕迹；要判断基层是否照办，仍须寻找县档、账簿、契约或后续案件。账本所示的研究方向是财政账册、市场网络、灾害环境与区域经济研究。可核的条件包括盐课是重要财政来源，既有商人网络与官府规制交错。应分层观察的后果是制度命令不能直接推知盐价、私盐或商人收益。证据界限：以乾隆十八年（1753年）的同期文书为定位起点；官修记录、奏报或外交文本服务特定机构立场，具体日期、规模、地方执行与对手视角须以独立材料复核。
\hypertarget{social:EQ-1754-AMURSANA-DEFECTION}{}\paragraph{乾隆十九年（1754年）}阿睦尔撒纳等准噶尔贵族向清廷归附，改变西北战争的联盟格局。从地方治理、身份、家庭与物质生活的角度，这一节点要求追问它在何处改变了资源、道路、组织或日常风险。核心来源为《清高宗实录》乾隆十九年相关条；军机处档、理藩院档或相关方略。这类编年材料可以标定朝廷获知和叙述事件的时间，却往往压低妇女、儿童、佃户、流民和非文书精英的行动。账本所示的研究方向是多语边疆档案、盟旗/寺院/绿洲地方史与帝国治理研究。可核的条件包括准噶尔内部权力竞争加深。应分层观察的后果是清廷将其视为进军伊犁的机会，后续关系迅速破裂。证据界限：以乾隆十九年（1754年）的同期文书为定位起点；官修记录、奏报或外交文本服务特定机构立场，具体日期、规模、地方执行与对手视角须以独立材料复核。
\hypertarget{social:EQ-1755-ILI-EXPEDITION}{}\paragraph{乾隆二十年（1755年）}清军分路进入伊犁，击溃达瓦齐政权并占领准噶尔核心地区。从地方治理、身份、家庭与物质生活的角度，这一节点要求追问它在何处改变了资源、道路、组织或日常风险。核心来源为《清高宗实录》乾隆二十年相关条；军机处档、战事方略及地方奏报。这类编年材料可以标定朝廷获知和叙述事件的时间，却往往压低妇女、儿童、佃户、流民和非文书精英的行动。账本所示的研究方向是战争动员、军需、战场暴力与地方社会研究。可核的条件包括准噶尔内部裂解和清军长期准备相互配合。应分层观察的后果是占领并未立即稳定，阿睦尔撒纳再起使战争延续。证据界限：以乾隆二十年（1755年）的同期文书为定位起点；官修记录、奏报或外交文本服务特定机构立场，具体日期、规模、地方执行与对手视角须以独立材料复核。
\hypertarget{social:EQ-1755-AMURSANA-REVOLT}{}\paragraph{乾隆二十年（1755年）}阿睦尔撒纳反清，准噶尔战后秩序重新陷入军事冲突。从地方治理、身份、家庭与物质生活的角度，这一节点要求追问它在何处改变了资源、道路、组织或日常风险。核心来源为《清高宗实录》乾隆二十年相关条；军机处档、战事方略及地方奏报。这类编年材料可以标定朝廷获知和叙述事件的时间，却往往压低妇女、儿童、佃户、流民和非文书精英的行动。账本所示的研究方向是战争动员、军需、战场暴力与地方社会研究。可核的条件包括清廷对新归附力量的安置与阿睦尔撒纳的权力期待冲突。应分层观察的后果是清军扩大报复性军事行动，人口损失叙述须谨慎核验。证据界限：以乾隆二十年（1755年）的同期文书为定位起点；官修记录、奏报或外交文本服务特定机构立场，具体日期、规模、地方执行与对手视角须以独立材料复核。
\hypertarget{social:EQ-1756-BANNER-CIVILIAN-REGISTRATION}{}\paragraph{乾隆二十一年（1756年）}清廷命令部分八旗副户口改归民籍，继续缓解旗籍供养压力。从地方治理、身份、家庭与物质生活的角度，这一节点要求追问它在何处改变了资源、道路、组织或日常风险。核心来源为《清高宗实录》乾隆二十一年相关条；地方志、督抚奏折及地方档案。这类编年材料可以标定朝廷获知和叙述事件的时间，却往往压低妇女、儿童、佃户、流民和非文书精英的行动。账本所示的研究方向是地方档案、迁徙、宗教、族群与社会动员研究。可核的条件包括旗内人口与财政资源不匹配。应分层观察的后果是身份转变在地方具有不同的程序和实际结果。证据界限：以乾隆二十一年（1756年）的同期文书为定位起点；官修记录、奏报或外交文本服务特定机构立场，具体日期、规模、地方执行与对手视角须以独立材料复核。
\hypertarget{social:EQ-1757-CANTON-TRADE-RESTRICTION}{}\paragraph{乾隆二十二年（1757年）}清廷将西洋贸易集中于广州口岸，形成后来所谓“一口通商”的制度框架。从地方治理、身份、家庭与物质生活的角度，这一节点要求追问它在何处改变了资源、道路、组织或日常风险。核心来源为《清高宗实录》乾隆二十二年相关条；《筹办夷务始末》相关年卷与中外照会、条约文本。它首先留下的是官府分类、任命、处置或资源调拨的痕迹；要判断基层是否照办，仍须寻找县档、账簿、契约或后续案件。账本所示的研究方向是条约文本、外交档案、口岸社会与国际法实践研究。可核的条件包括沿海治安、关税和对外接触的管理需求上升。应分层观察的后果是广州贸易网络扩大，但禁限条文不等于所有海上往来消失。证据界限：以乾隆二十二年（1757年）的同期文书为定位起点；官修记录、奏报或外交文本服务特定机构立场，具体日期、规模、地方执行与对手视角须以独立材料复核。
\hypertarget{social:EQ-1757-AMURSANA-DEATH}{}\paragraph{乾隆二十二年（1757年）}阿睦尔撒纳逃入俄境后死亡，清准战争的关键对手退出。从地方治理、身份、家庭与物质生活的角度，这一节点要求追问它在何处改变了资源、道路、组织或日常风险。核心来源为《清高宗实录》乾隆二十二年相关条；军机处档、理藩院档或相关方略。这类编年材料可以标定朝廷获知和叙述事件的时间，却往往压低妇女、儿童、佃户、流民和非文书精英的行动。账本所示的研究方向是多语边疆档案、盟旗/寺院/绿洲地方史与帝国治理研究。可核的条件包括清军追击与俄清边境交涉交织。应分层观察的后果是其死亡与边境信息须结合俄方材料核对。证据界限：以乾隆二十二年（1757年）的同期文书为定位起点；官修记录、奏报或外交文本服务特定机构立场，具体日期、规模、地方执行与对手视角须以独立材料复核。
\hypertarget{social:EQ-1758-KHOJA-REBELLION}{}\paragraph{乾隆二十三年（1758年）}大小和卓反清，清军由北路向天山南路推进。从地方治理、身份、家庭与物质生活的角度，这一节点要求追问它在何处改变了资源、道路、组织或日常风险。核心来源为《清高宗实录》乾隆二十三年相关条；军机处档、战事方略及地方奏报。这类编年材料可以标定朝廷获知和叙述事件的时间，却往往压低妇女、儿童、佃户、流民和非文书精英的行动。账本所示的研究方向是战争动员、军需、战场暴力与地方社会研究。可核的条件包括伊犁占领后绿洲地区的权力安排尚未稳定。应分层观察的后果是清廷的军事胜利与绿洲社会的重新组织需分别追踪。证据界限：以乾隆二十三年（1758年）的同期文书为定位起点；官修记录、奏报或外交文本服务特定机构立场，具体日期、规模、地方执行与对手视角须以独立材料复核。
\hypertarget{social:EQ-1759-KASHGAR-YARKAND-CAMPAIGN}{}\paragraph{乾隆二十四年（1759年）}清军攻入喀什噶尔、叶尔羌等地，和卓政权覆亡。从地方治理、身份、家庭与物质生活的角度，这一节点要求追问它在何处改变了资源、道路、组织或日常风险。核心来源为《清高宗实录》乾隆二十四年相关条；军机处档、战事方略及地方奏报。这类编年材料可以标定朝廷获知和叙述事件的时间，却往往压低妇女、儿童、佃户、流民和非文书精英的行动。账本所示的研究方向是战争动员、军需、战场暴力与地方社会研究。可核的条件包括清军跨越天山并依赖绿洲补给和地方协力。应分层观察的后果是天山南路纳入清帝国的行政过程远晚于战役结束。证据界限：以乾隆二十四年（1759年）的同期文书为定位起点；官修记录、奏报或外交文本服务特定机构立场，具体日期、规模、地方执行与对手视角须以独立材料复核。
\subsection{1760—1769：地方材料所能看见的尺度}
以下事件按账本时间排列；它们不是同质的社会样本，而是提示应到何种地点、机构和材料中继续追问的入口。
\hypertarget{social:EQ-1760-ILI-GENERAL}{}\paragraph{乾隆二十五年（1760年）}清廷设伊犁将军并展开驻防、屯田与移民安置。从地方治理、身份、家庭与物质生活的角度，这一节点要求追问它在何处改变了资源、道路、组织或日常风险。核心来源为《清高宗实录》乾隆二十五年相关条；军机处档、理藩院档或相关方略。这类编年材料可以标定朝廷获知和叙述事件的时间，却往往压低妇女、儿童、佃户、流民和非文书精英的行动。账本所示的研究方向是多语边疆档案、盟旗/寺院/绿洲地方史与帝国治理研究。可核的条件包括西北战后需要长期军政管理和粮食供给。应分层观察的后果是驻防城市、屯田和地方社会间的关系不可简化为单向控制。证据界限：以乾隆二十五年（1760年）的同期文书为定位起点；官修记录、奏报或外交文本服务特定机构立场，具体日期、规模、地方执行与对手视角须以独立材料复核。
\hypertarget{social:EQ-1761-XINJIANG-GARRISONS}{}\paragraph{乾隆二十六年（1761年）}清廷在新疆配置满汉蒙军队和绿营驻防，建立军政节点。从地方治理、身份、家庭与物质生活的角度，这一节点要求追问它在何处改变了资源、道路、组织或日常风险。核心来源为《清高宗实录》乾隆二十六年相关条；军机处档、理藩院档或相关方略。这类编年材料可以标定朝廷获知和叙述事件的时间，却往往压低妇女、儿童、佃户、流民和非文书精英的行动。账本所示的研究方向是多语边疆档案、盟旗/寺院/绿洲地方史与帝国治理研究。可核的条件包括新征服地区交通遥远且人口与资源分布不均。应分层观察的后果是兵额与驻军实际存在需以档案、城址和账册互证。证据界限：以乾隆二十六年（1761年）的同期文书为定位起点；官修记录、奏报或外交文本服务特定机构立场，具体日期、规模、地方执行与对手视角须以独立材料复核。
\hypertarget{social:EQ-1762-URUMQI-FOUNDATION}{}\paragraph{乾隆二十七年（1762年）}迪化等北疆城市的军政建设加快，成为屯田和交通枢纽。从地方治理、身份、家庭与物质生活的角度，这一节点要求追问它在何处改变了资源、道路、组织或日常风险。核心来源为《清高宗实录》乾隆二十七年相关条；军机处档、理藩院档或相关方略。这类编年材料可以标定朝廷获知和叙述事件的时间，却往往压低妇女、儿童、佃户、流民和非文书精英的行动。账本所示的研究方向是多语边疆档案、盟旗/寺院/绿洲地方史与帝国治理研究。可核的条件包括伊犁—东疆之间需要稳定的补给与行政据点。应分层观察的后果是城市建设改变土地与水利使用，地方社会后果需另证。证据界限：以乾隆二十七年（1762年）的同期文书为定位起点；官修记录、奏报或外交文本服务特定机构立场，具体日期、规模、地方执行与对手视角须以独立材料复核。
\hypertarget{social:EQ-1763-XINJIANG-SETTLEMENT}{}\paragraph{乾隆二十八年（1763年）}清廷继续在新疆安置屯田人口、军户和商业网络。从地方治理、身份、家庭与物质生活的角度，这一节点要求追问它在何处改变了资源、道路、组织或日常风险。核心来源为《清高宗实录》乾隆二十八年相关条；地方志、督抚奏折及地方档案。这类编年材料可以标定朝廷获知和叙述事件的时间，却往往压低妇女、儿童、佃户、流民和非文书精英的行动。账本所示的研究方向是地方档案、迁徙、宗教、族群与社会动员研究。可核的条件包括驻防粮食、劳力和市场供给依赖移民与地方生产。应分层观察的后果是户籍与屯田册只能说明登记，不自动等于稳定定居。证据界限：以乾隆二十八年（1763年）的同期文书为定位起点；官修记录、奏报或外交文本服务特定机构立场，具体日期、规模、地方执行与对手视角须以独立材料复核。
\hypertarget{social:EQ-1764-MANCHU-BANNER-RESETTLEMENT}{}\paragraph{乾隆二十九年（1764年）}部分八旗兵丁被调往西北驻防，旗籍家庭的迁移与供给成为边务问题。从地方治理、身份、家庭与物质生活的角度，这一节点要求追问它在何处改变了资源、道路、组织或日常风险。核心来源为《清高宗实录》乾隆二十九年相关条；地方志、督抚奏折及地方档案。这类编年材料可以标定朝廷获知和叙述事件的时间，却往往压低妇女、儿童、佃户、流民和非文书精英的行动。账本所示的研究方向是地方档案、迁徙、宗教、族群与社会动员研究。可核的条件包括伊犁和北疆需要长期守备力量。应分层观察的后果是迁徙改变旗人家庭生活，不能只以军政文书描述。证据界限：以乾隆二十九年（1764年）的同期文书为定位起点；官修记录、奏报或外交文本服务特定机构立场，具体日期、规模、地方执行与对手视角须以独立材料复核。
\hypertarget{social:EQ-1765-USH-REBELLION}{}\paragraph{乾隆三十年（1765年）}乌什爆发反清事件，清军镇压并重组当地治理。从地方治理、身份、家庭与物质生活的角度，这一节点要求追问它在何处改变了资源、道路、组织或日常风险。核心来源为《清高宗实录》乾隆三十年相关条；军机处档、战事方略及地方奏报。这类编年材料可以标定朝廷获知和叙述事件的时间，却往往压低妇女、儿童、佃户、流民和非文书精英的行动。账本所示的研究方向是战争动员、军需、战场暴力与地方社会研究。可核的条件包括地方征发、驻军关系和绿洲权力矛盾累积。应分层观察的后果是官方定性和伤亡数字须与维吾尔文及地方材料交叉检验。证据界限：以乾隆三十年（1765年）的同期文书为定位起点；官修记录、奏报或外交文本服务特定机构立场，具体日期、规模、地方执行与对手视角须以独立材料复核。
\hypertarget{social:EQ-1766-BURMA-WAR-BEGIN}{}\paragraph{乾隆三十一年（1766年）}清军对缅甸发动远征，西南边境战争展开。从地方治理、身份、家庭与物质生活的角度，这一节点要求追问它在何处改变了资源、道路、组织或日常风险。核心来源为《清高宗实录》乾隆三十一年相关条；军机处档、战事方略及地方奏报。这类编年材料可以标定朝廷获知和叙述事件的时间，却往往压低妇女、儿童、佃户、流民和非文书精英的行动。账本所示的研究方向是战争动员、军需、战场暴力与地方社会研究。可核的条件包括边境贸易、土司关系和清缅政治冲突相互交织。应分层观察的后果是山地疾病、补给和气候削弱远征军的战斗力。证据界限：以乾隆三十一年（1766年）的同期文书为定位起点；官修记录、奏报或外交文本服务特定机构立场，具体日期、规模、地方执行与对手视角须以独立材料复核。
\hypertarget{social:EQ-1767-BURMA-CAMPAIGN}{}\paragraph{乾隆三十二年（1767年）}清缅战争持续，清军在边境和山地路线中遭遇重大困难。从地方治理、身份、家庭与物质生活的角度，这一节点要求追问它在何处改变了资源、道路、组织或日常风险。核心来源为《清高宗实录》乾隆三十二年相关条；军机处档、战事方略及地方奏报。这类编年材料可以标定朝廷获知和叙述事件的时间，却往往压低妇女、儿童、佃户、流民和非文书精英的行动。账本所示的研究方向是战争动员、军需、战场暴力与地方社会研究。可核的条件包括跨区域指挥与热带环境限制北方军队行动。应分层观察的后果是军事文书的战果叙述须与缅方材料并读。证据界限：以乾隆三十二年（1767年）的同期文书为定位起点；官修记录、奏报或外交文本服务特定机构立场，具体日期、规模、地方执行与对手视角须以独立材料复核。
\hypertarget{social:EQ-1768-LITERARY-INQUISITION}{}\paragraph{乾隆三十三年（1768年）}乾隆朝以文字狱案件强化对著述、地方士人和政治语言的控制。从地方治理、身份、家庭与物质生活的角度，这一节点要求追问它在何处改变了资源、道路、组织或日常风险。核心来源为《清高宗实录》乾隆三十三年相关条；内阁档、文集或报刊相关记录。这类编年材料可以标定朝廷获知和叙述事件的时间，却往往压低妇女、儿童、佃户、流民和非文书精英的行动。账本所示的研究方向是知识生产、教育、宗教与公共传播研究。可核的条件包括朝廷担忧地方文人以历史、族群或政论表达异议。应分层观察的后果是案件卷宗显示国家分类，不代表所有读书人的意图。证据界限：以乾隆三十三年（1768年）的同期文书为定位起点；官修记录、奏报或外交文本服务特定机构立场，具体日期、规模、地方执行与对手视角须以独立材料复核。
\hypertarget{social:EQ-1768-SECOND-JINCHUAN-WAR}{}\paragraph{乾隆三十三年（1768年）}第二次金川战争开始，清廷投入更大兵力和资源。从地方治理、身份、家庭与物质生活的角度，这一节点要求追问它在何处改变了资源、道路、组织或日常风险。核心来源为《清高宗实录》乾隆三十三年相关条；军机处档、战事方略及地方奏报。这类编年材料可以标定朝廷获知和叙述事件的时间，却往往压低妇女、儿童、佃户、流民和非文书精英的行动。账本所示的研究方向是战争动员、军需、战场暴力与地方社会研究。可核的条件包括第一次和议未能解决土司与清廷的权力冲突。应分层观察的后果是长期山地战带来巨额军费与地方劳役负担。证据界限：以乾隆三十三年（1768年）的同期文书为定位起点；官修记录、奏报或外交文本服务特定机构立场，具体日期、规模、地方执行与对手视角须以独立材料复核。
\hypertarget{social:EQ-1769-BURMA-SETTLEMENT}{}\paragraph{乾隆三十四年（1769年）}清缅战争在反复交战后趋向停战和使节往来安排。从地方治理、身份、家庭与物质生活的角度，这一节点要求追问它在何处改变了资源、道路、组织或日常风险。核心来源为《清高宗实录》乾隆三十四年相关条；《筹办夷务始末》相关年卷与中外照会、条约文本。它首先留下的是官府分类、任命、处置或资源调拨的痕迹；要判断基层是否照办，仍须寻找县档、账簿、契约或后续案件。账本所示的研究方向是条约文本、外交档案、口岸社会与国际法实践研究。可核的条件包括双方均承受远征成本与边境治理压力。应分层观察的后果是朝贡或使节语言不能直接等同单一的宗主权关系。证据界限：以乾隆三十四年（1769年）的同期文书为定位起点；官修记录、奏报或外交文本服务特定机构立场，具体日期、规模、地方执行与对手视角须以独立材料复核。
\hypertarget{social:EQ-1769-JINCHUAN-ESCALATION}{}\paragraph{乾隆三十四年（1769年）}金川战场的道路、炮兵、粮运与将领调度不断升级。从地方治理、身份、家庭与物质生活的角度，这一节点要求追问它在何处改变了资源、道路、组织或日常风险。核心来源为《清高宗实录》乾隆三十四年相关条；军机处档、战事方略及地方奏报。这类编年材料可以标定朝廷获知和叙述事件的时间，却往往压低妇女、儿童、佃户、流民和非文书精英的行动。账本所示的研究方向是战争动员、军需、战场暴力与地方社会研究。可核的条件包括山地堡寨体系使常规野战难以奏效。应分层观察的后果是清廷战后更重视边地交通与驻防，但其效果需地方材料检验。证据界限：以乾隆三十四年（1769年）的同期文书为定位起点；官修记录、奏报或外交文本服务特定机构立场，具体日期、规模、地方执行与对手视角须以独立材料复核。
\subsection{1770—1779：地方材料所能看见的尺度}
以下事件按账本时间排列；它们不是同质的社会样本，而是提示应到何种地点、机构和材料中继续追问的入口。
\hypertarget{social:EQ-1770-SICHUAN-SUPPLY-BURDEN}{}\paragraph{乾隆三十五年（1770年）}金川军需加重四川及邻省的转运、民夫和财政负担。从地方治理、身份、家庭与物质生活的角度，这一节点要求追问它在何处改变了资源、道路、组织或日常风险。核心来源为《清高宗实录》乾隆三十五年相关条；户部奏销、盐政、河工或海关档案。它首先留下的是官府分类、任命、处置或资源调拨的痕迹；要判断基层是否照办，仍须寻找县档、账簿、契约或后续案件。账本所示的研究方向是财政账册、市场网络、灾害环境与区域经济研究。可核的条件包括大规模驻军消耗粮食、牲畜和道路资源。应分层观察的后果是战时征发对村社和市场的影响不应被军功叙事遮蔽。证据界限：以乾隆三十五年（1770年）的同期文书为定位起点；官修记录、奏报或外交文本服务特定机构立场，具体日期、规模、地方执行与对手视角须以独立材料复核。
\hypertarget{social:EQ-1771-TORGHUT-RETURN}{}\paragraph{乾隆三十六年（1771年）}土尔扈特部自伏尔加回归，清廷在伊犁等地安置部众。从地方治理、身份、家庭与物质生活的角度，这一节点要求追问它在何处改变了资源、道路、组织或日常风险。核心来源为《清高宗实录》乾隆三十六年相关条；军机处档、理藩院档或相关方略。这类编年材料可以标定朝廷获知和叙述事件的时间，却往往压低妇女、儿童、佃户、流民和非文书精英的行动。账本所示的研究方向是多语边疆档案、盟旗/寺院/绿洲地方史与帝国治理研究。可核的条件包括俄国扩张、草原压力与部众迁徙决策共同作用。应分层观察的后果是安置文书不能掩盖迁徙损失和内部差异。证据界限：以乾隆三十六年（1771年）的同期文书为定位起点；官修记录、奏报或外交文本服务特定机构立场，具体日期、规模、地方执行与对手视角须以独立材料复核。
\hypertarget{social:EQ-1772-SIKU-QUANSHU-ORDER}{}\paragraph{乾隆三十七年（1772年）}乾隆帝下令征集、编纂《四库全书》，建立全国性的图书征集机制。从地方治理、身份、家庭与物质生活的角度，这一节点要求追问它在何处改变了资源、道路、组织或日常风险。核心来源为《清高宗实录》乾隆三十七年相关条；内阁档、文集或报刊相关记录。这类编年材料可以标定朝廷获知和叙述事件的时间，却往往压低妇女、儿童、佃户、流民和非文书精英的行动。账本所示的研究方向是知识生产、教育、宗教与公共传播研究。可核的条件包括朝廷希望整理典籍并审查政治敏感著作。应分层观察的后果是征书和禁毁并行，书目不能等于实际保存或流通。证据界限：以乾隆三十七年（1772年）的同期文书为定位起点；官修记录、奏报或外交文本服务特定机构立场，具体日期、规模、地方执行与对手视角须以独立材料复核。
\hypertarget{social:EQ-1773-SICHUAN-RECRUITMENT}{}\paragraph{乾隆三十八年（1773年）}金川战事促使清廷继续在四川及邻省征调兵丁和夫役。从地方治理、身份、家庭与物质生活的角度，这一节点要求追问它在何处改变了资源、道路、组织或日常风险。核心来源为《清高宗实录》乾隆三十八年相关条；军机处档、战事方略及地方奏报。这类编年材料可以标定朝廷获知和叙述事件的时间，却往往压低妇女、儿童、佃户、流民和非文书精英的行动。账本所示的研究方向是战争动员、军需、战场暴力与地方社会研究。可核的条件包括战场消耗超过初期预期。应分层观察的后果是兵役与夫役的地方分配存在显著不均。证据界限：以乾隆三十八年（1773年）的同期文书为定位起点；官修记录、奏报或外交文本服务特定机构立场，具体日期、规模、地方执行与对手视角须以独立材料复核。
\hypertarget{social:EQ-1774-WANG-LUN-UPRISING}{}\paragraph{乾隆三十九年（1774年）}王伦在山东起事，清廷迅速调兵镇压。从地方治理、身份、家庭与物质生活的角度，这一节点要求追问它在何处改变了资源、道路、组织或日常风险。核心来源为《清高宗实录》乾隆三十九年相关条；地方志、督抚奏折及地方档案。这类编年材料可以标定朝廷获知和叙述事件的时间，却往往压低妇女、儿童、佃户、流民和非文书精英的行动。账本所示的研究方向是地方档案、迁徙、宗教、族群与社会动员研究。可核的条件包括宗教结社、地方贫困和治安压力形成起事背景。应分层观察的后果是官府对“教门”的分类不应取代参与者多样诉求。证据界限：以乾隆三十九年（1774年）的同期文书为定位起点；官修记录、奏报或外交文本服务特定机构立场，具体日期、规模、地方执行与对手视角须以独立材料复核。
\hypertarget{social:EQ-1775-SIKU-COLLECTION}{}\paragraph{乾隆四十年（1775年）}《四库全书》征书与审查在各省展开，地方藏书和文人网络受到影响。从地方治理、身份、家庭与物质生活的角度，这一节点要求追问它在何处改变了资源、道路、组织或日常风险。核心来源为《清高宗实录》乾隆四十年相关条；内阁档、文集或报刊相关记录。这类编年材料可以标定朝廷获知和叙述事件的时间，却往往压低妇女、儿童、佃户、流民和非文书精英的行动。账本所示的研究方向是知识生产、教育、宗教与公共传播研究。可核的条件包括中央需要将地方文献纳入可控的知识秩序。应分层观察的后果是禁毁目录和地方执行范围应与存书、抄本和书信互证。证据界限：以乾隆四十年（1775年）的同期文书为定位起点；官修记录、奏报或外交文本服务特定机构立场，具体日期、规模、地方执行与对手视角须以独立材料复核。
\hypertarget{social:EQ-1776-JINCHUAN-CONQUEST}{}\paragraph{乾隆四十一年（1776年）}清军攻克金川，第二次金川战争结束。从地方治理、身份、家庭与物质生活的角度，这一节点要求追问它在何处改变了资源、道路、组织或日常风险。核心来源为《清高宗实录》乾隆四十一年相关条；军机处档、战事方略及地方奏报。这类编年材料可以标定朝廷获知和叙述事件的时间，却往往压低妇女、儿童、佃户、流民和非文书精英的行动。账本所示的研究方向是战争动员、军需、战场暴力与地方社会研究。可核的条件包括长期工程、火器和大规模转运最终压倒山地防御。应分层观察的后果是战后改土归流与人口迁移的过程不应由“平定”一词遮蔽。证据界限：以乾隆四十一年（1776年）的同期文书为定位起点；官修记录、奏报或外交文本服务特定机构立场，具体日期、规模、地方执行与对手视角须以独立材料复核。
\hypertarget{social:EQ-1777-JINCHUAN-ADMINISTRATIVE-RESET}{}\paragraph{乾隆四十二年（1777年）}清廷在金川地区推行改土归流、设治和驻防安排。从地方治理、身份、家庭与物质生活的角度，这一节点要求追问它在何处改变了资源、道路、组织或日常风险。核心来源为《清高宗实录》乾隆四十二年相关条；军机处档、理藩院档或相关方略。这类编年材料可以标定朝廷获知和叙述事件的时间，却往往压低妇女、儿童、佃户、流民和非文书精英的行动。账本所示的研究方向是多语边疆档案、盟旗/寺院/绿洲地方史与帝国治理研究。可核的条件包括军事胜利后需要将战果转为持续行政控制。应分层观察的后果是地方秩序重组伴随土地、劳役和宗教网络变化。证据界限：以乾隆四十二年（1777年）的同期文书为定位起点；官修记录、奏报或外交文本服务特定机构立场，具体日期、规模、地方执行与对手视角须以独立材料复核。
\hypertarget{social:EQ-1778-SICHUAN-TAX-RECOVERY}{}\paragraph{乾隆四十三年（1778年）}金川战后四川财政和民力恢复成为地方官奏报的重要议题。从地方治理、身份、家庭与物质生活的角度，这一节点要求追问它在何处改变了资源、道路、组织或日常风险。核心来源为《清高宗实录》乾隆四十三年相关条；户部奏销、盐政、河工或海关档案。它首先留下的是官府分类、任命、处置或资源调拨的痕迹；要判断基层是否照办，仍须寻找县档、账簿、契约或后续案件。账本所示的研究方向是财政账册、市场网络、灾害环境与区域经济研究。可核的条件包括长期战争造成钱粮、道路和人口损耗。应分层观察的后果是官方恢复叙事须同地方账册和市场材料核对。证据界限：以乾隆四十三年（1778年）的同期文书为定位起点；官修记录、奏报或外交文本服务特定机构立场，具体日期、规模、地方执行与对手视角须以独立材料复核。
\hypertarget{social:EQ-1779-VIETNAM-INTERVENTION}{}\paragraph{乾隆四十四年（1779年）}清廷介入安南政局并派兵南下，孙士毅军一度进入升龙。从地方治理、身份、家庭与物质生活的角度，这一节点要求追问它在何处改变了资源、道路、组织或日常风险。核心来源为《清高宗实录》乾隆四十四年相关条；军机处档、战事方略及地方奏报。这类编年材料可以标定朝廷获知和叙述事件的时间，却往往压低妇女、儿童、佃户、流民和非文书精英的行动。账本所示的研究方向是战争动员、军需、战场暴力与地方社会研究。可核的条件包括黎朝复辟请求与区域权力竞争相交。应分层观察的后果是清军行动受当地政治与补给限制，战果叙述有待中越材料互校。证据界限：以乾隆四十四年（1779年）的同期文书为定位起点；官修记录、奏报或外交文本服务特定机构立场，具体日期、规模、地方执行与对手视角须以独立材料复核。
\subsection{1780—1789：地方材料所能看见的尺度}
以下事件按账本时间排列；它们不是同质的社会样本，而是提示应到何种地点、机构和材料中继续追问的入口。
\hypertarget{social:EQ-1780-VIETNAM-WITHDRAWAL}{}\paragraph{乾隆四十五年（1780年）}清军撤离安南，双方通过使节和册封语言恢复关系。从地方治理、身份、家庭与物质生活的角度，这一节点要求追问它在何处改变了资源、道路、组织或日常风险。核心来源为《清高宗实录》乾隆四十五年相关条；《筹办夷务始末》相关年卷与中外照会、条约文本。它首先留下的是官府分类、任命、处置或资源调拨的痕迹；要判断基层是否照办，仍须寻找县档、账簿、契约或后续案件。账本所示的研究方向是条约文本、外交档案、口岸社会与国际法实践研究。可核的条件包括军事行动未能建立稳定的复辟政权。应分层观察的后果是外交礼仪提供缓和机制，不等于消除区域竞争。证据界限：以乾隆四十五年（1780年）的同期文书为定位起点；官修记录、奏报或外交文本服务特定机构立场，具体日期、规模、地方执行与对手视角须以独立材料复核。
\hypertarget{social:EQ-1781-GANSU-MUSLIM-UPRISING}{}\paragraph{乾隆四十六年（1781年）}甘肃爆发回民起事，清军实施镇压与后续治理。从地方治理、身份、家庭与物质生活的角度，这一节点要求追问它在何处改变了资源、道路、组织或日常风险。核心来源为《清高宗实录》乾隆四十六年相关条；军机处档、战事方略及地方奏报。这类编年材料可以标定朝廷获知和叙述事件的时间，却往往压低妇女、儿童、佃户、流民和非文书精英的行动。账本所示的研究方向是战争动员、军需、战场暴力与地方社会研究。可核的条件包括地方纠纷、行政失当和族群分类被动员为冲突语言。应分层观察的后果是伤亡、动机与群体边界必须以地方、多语材料审慎处理。证据界限：以乾隆四十六年（1781年）的同期文书为定位起点；官修记录、奏报或外交文本服务特定机构立场，具体日期、规模、地方执行与对手视角须以独立材料复核。
\hypertarget{social:EQ-1782-GANSU-POSTWAR-GOVERNANCE}{}\paragraph{乾隆四十七年（1782年）}清廷在甘肃起事后调整驻兵、司法和地方控制措施。从地方治理、身份、家庭与物质生活的角度，这一节点要求追问它在何处改变了资源、道路、组织或日常风险。核心来源为《清高宗实录》乾隆四十七年相关条；地方志、督抚奏折及地方档案。这类编年材料可以标定朝廷获知和叙述事件的时间，却往往压低妇女、儿童、佃户、流民和非文书精英的行动。账本所示的研究方向是地方档案、迁徙、宗教、族群与社会动员研究。可核的条件包括战后需要恢复交通、赋役与社区关系。应分层观察的后果是官方治安叙事不能替代地方生计和迁徙记录。证据界限：以乾隆四十七年（1782年）的同期文书为定位起点；官修记录、奏报或外交文本服务特定机构立场，具体日期、规模、地方执行与对手视角须以独立材料复核。
\hypertarget{social:EQ-1783-SALT-MONOPOLY-REVIEW}{}\paragraph{乾隆四十八年（1783年）}清廷继续讨论盐政弊端、商人负担和课额征解。从地方治理、身份、家庭与物质生活的角度，这一节点要求追问它在何处改变了资源、道路、组织或日常风险。核心来源为《清高宗实录》乾隆四十八年相关条；户部奏销、盐政、河工或海关档案。它首先留下的是官府分类、任命、处置或资源调拨的痕迹；要判断基层是否照办，仍须寻找县档、账簿、契约或后续案件。账本所示的研究方向是财政账册、市场网络、灾害环境与区域经济研究。可核的条件包括盐课、运销网络与地方财政紧密相连。应分层观察的后果是改革方案须与实际盐价、私盐和账册区分。证据界限：以乾隆四十八年（1783年）的同期文书为定位起点；官修记录、奏报或外交文本服务特定机构立场，具体日期、规模、地方执行与对手视角须以独立材料复核。
\hypertarget{social:EQ-1785-TAIWAN-GARRISON}{}\paragraph{乾隆五十年（1785年）}清廷调整台湾驻防、渡台与地方治安安排。从地方治理、身份、家庭与物质生活的角度，这一节点要求追问它在何处改变了资源、道路、组织或日常风险。核心来源为《清高宗实录》乾隆五十年相关条；军机处档、理藩院档或相关方略。这类编年材料可以标定朝廷获知和叙述事件的时间，却往往压低妇女、儿童、佃户、流民和非文书精英的行动。账本所示的研究方向是多语边疆档案、盟旗/寺院/绿洲地方史与帝国治理研究。可核的条件包括海峡调兵和岛内治理成本持续上升。应分层观察的后果是制度设计与港口、村社的实际执行需要分开考察。证据界限：以乾隆五十年（1785年）的同期文书为定位起点；官修记录、奏报或外交文本服务特定机构立场，具体日期、规模、地方执行与对手视角须以独立材料复核。
\hypertarget{social:EQ-1786-LIN-SHUANGWEN-UPRISING}{}\paragraph{乾隆五十一年（1786年）}林爽文起事在台湾爆发并迅速扩展。从地方治理、身份、家庭与物质生活的角度，这一节点要求追问它在何处改变了资源、道路、组织或日常风险。核心来源为《清高宗实录》乾隆五十一年相关条；军机处档、战事方略及地方奏报。这类编年材料可以标定朝廷获知和叙述事件的时间，却往往压低妇女、儿童、佃户、流民和非文书精英的行动。账本所示的研究方向是战争动员、军需、战场暴力与地方社会研究。可核的条件包括地方权力竞争、赋役与治安危机叠加。应分层观察的后果是清廷须跨海调兵，地方参与者的目标并不一致。证据界限：以乾隆五十一年（1786年）的同期文书为定位起点；官修记录、奏报或外交文本服务特定机构立场，具体日期、规模、地方执行与对手视角须以独立材料复核。
\hypertarget{social:EQ-1786-FIRST-GURKHA-WAR}{}\paragraph{乾隆五十一年（1786年）}廓尔喀军进入西藏边境，清廷开始处理高原军事危机。从地方治理、身份、家庭与物质生活的角度，这一节点要求追问它在何处改变了资源、道路、组织或日常风险。核心来源为《清高宗实录》乾隆五十一年相关条；军机处档、理藩院档或相关方略。这类编年材料可以标定朝廷获知和叙述事件的时间，却往往压低妇女、儿童、佃户、流民和非文书精英的行动。账本所示的研究方向是多语边疆档案、盟旗/寺院/绿洲地方史与帝国治理研究。可核的条件包括贸易、货币与边界纠纷恶化。应分层观察的后果是清藏与廓尔喀材料对冲突起因有不同叙述。证据界限：以乾隆五十一年（1786年）的同期文书为定位起点；官修记录、奏报或外交文本服务特定机构立场，具体日期、规模、地方执行与对手视角须以独立材料复核。
\hypertarget{social:EQ-1787-TAIWAN-CAMPAIGN}{}\paragraph{乾隆五十二年（1787年）}福康安等率军赴台镇压林爽文起事，海峡运输成为关键。从地方治理、身份、家庭与物质生活的角度，这一节点要求追问它在何处改变了资源、道路、组织或日常风险。核心来源为《清高宗实录》乾隆五十二年相关条；军机处档、战事方略及地方奏报。这类编年材料可以标定朝廷获知和叙述事件的时间，却往往压低妇女、儿童、佃户、流民和非文书精英的行动。账本所示的研究方向是战争动员、军需、战场暴力与地方社会研究。可核的条件包括岛内守军不足且起事波及多地。应分层观察的后果是军事镇压与地方招抚、清算应分别记录。证据界限：以乾隆五十二年（1787年）的同期文书为定位起点；官修记录、奏报或外交文本服务特定机构立场，具体日期、规模、地方执行与对手视角须以独立材料复核。
\hypertarget{social:EQ-1788-LIN-SHUANGWEN-DEFEAT}{}\paragraph{乾隆五十三年（1788年）}清军击败林爽文，台湾起事基本结束。从地方治理、身份、家庭与物质生活的角度，这一节点要求追问它在何处改变了资源、道路、组织或日常风险。核心来源为《清高宗实录》乾隆五十三年相关条；军机处档、战事方略及地方奏报。这类编年材料可以标定朝廷获知和叙述事件的时间，却往往压低妇女、儿童、佃户、流民和非文书精英的行动。账本所示的研究方向是战争动员、军需、战场暴力与地方社会研究。可核的条件包括增援军队、地方协力和起事内部困难共同影响战局。应分层观察的后果是战后处置、伤亡和地方秩序重建不能由官书凯旋叙事概括。证据界限：以乾隆五十三年（1788年）的同期文书为定位起点；官修记录、奏报或外交文本服务特定机构立场，具体日期、规模、地方执行与对手视角须以独立材料复核。
\hypertarget{social:EQ-1788-SECOND-GURKHA-WAR}{}\paragraph{乾隆五十三年（1788年）}廓尔喀再次入侵西藏，清军组织更大规模高原远征。从地方治理、身份、家庭与物质生活的角度，这一节点要求追问它在何处改变了资源、道路、组织或日常风险。核心来源为《清高宗实录》乾隆五十三年相关条；军机处档、战事方略及地方奏报。这类编年材料可以标定朝廷获知和叙述事件的时间，却往往压低妇女、儿童、佃户、流民和非文书精英的行动。账本所示的研究方向是战争动员、军需、战场暴力与地方社会研究。可核的条件包括前次安排未化解边境贸易与政治矛盾。应分层观察的后果是军需和高原路线使战争超出拉萨单一政治中心。证据界限：以乾隆五十三年（1788年）的同期文书为定位起点；官修记录、奏报或外交文本服务特定机构立场，具体日期、规模、地方执行与对手视角须以独立材料复核。
\hypertarget{social:EQ-1789-TIBET-NEPAL-CAMPAIGN}{}\paragraph{乾隆五十四年（1789年）}清军越过喜马拉雅方向追击廓尔喀，边境协商与军事行动并行。从地方治理、身份、家庭与物质生活的角度，这一节点要求追问它在何处改变了资源、道路、组织或日常风险。核心来源为《清高宗实录》乾隆五十四年相关条；军机处档、战事方略及地方奏报。这类编年材料可以标定朝廷获知和叙述事件的时间，却往往压低妇女、儿童、佃户、流民和非文书精英的行动。账本所示的研究方向是战争动员、军需、战场暴力与地方社会研究。可核的条件包括清廷意在重建西藏边防威慑。应分层观察的后果是行军路线、战果和伤亡须以藏、尼泊尔及清方材料核验。证据界限：以乾隆五十四年（1789年）的同期文书为定位起点；官修记录、奏报或外交文本服务特定机构立场，具体日期、规模、地方执行与对手视角须以独立材料复核。
\subsection{1790—1799：地方材料所能看见的尺度}
以下事件按账本时间排列；它们不是同质的社会样本，而是提示应到何种地点、机构和材料中继续追问的入口。
\hypertarget{social:EQ-1790-IMPERIAL-REGULATIONS-TIBET}{}\paragraph{乾隆五十五年（1790年）}清廷制定整顿西藏事务的制度性安排，强化驻藏大臣和金瓶掣签等机制。从地方治理、身份、家庭与物质生活的角度，这一节点要求追问它在何处改变了资源、道路、组织或日常风险。核心来源为《清高宗实录》乾隆五十五年相关条；军机处档、理藩院档或相关方略。这类编年材料可以标定朝廷获知和叙述事件的时间，却往往压低妇女、儿童、佃户、流民和非文书精英的行动。账本所示的研究方向是多语边疆档案、盟旗/寺院/绿洲地方史与帝国治理研究。可核的条件包括廓尔喀战争暴露既有治理安排的漏洞。应分层观察的后果是制度条文与其后宗教、地方政治实践不能直接画等号。证据界限：以乾隆五十五年（1790年）的同期文书为定位起点；官修记录、奏报或外交文本服务特定机构立场，具体日期、规模、地方执行与对手视角须以独立材料复核。
\hypertarget{social:EQ-1791-WHITE-LOTUS-PRECURSORS}{}\paragraph{乾隆五十六年（1791年）}川楚陕基层社会的教门网络和地方压力成为白莲教战争的背景。从地方治理、身份、家庭与物质生活的角度，这一节点要求追问它在何处改变了资源、道路、组织或日常风险。核心来源为《清高宗实录》乾隆五十六年相关条；地方志、督抚奏折及地方档案。这类编年材料可以标定朝廷获知和叙述事件的时间，却往往压低妇女、儿童、佃户、流民和非文书精英的行动。账本所示的研究方向是地方档案、迁徙、宗教、族群与社会动员研究。可核的条件包括人口、土地、赋役和地方治安矛盾长期积累。应分层观察的后果是官府分类为“邪教”的材料须与社区网络分开分析。证据界限：以乾隆五十六年（1791年）的同期文书为定位起点；官修记录、奏报或外交文本服务特定机构立场，具体日期、规模、地方执行与对手视角须以独立材料复核。
\hypertarget{social:EQ-1792-GURKHA-SETTLEMENT}{}\paragraph{乾隆五十七年（1792年）}清廷与廓尔喀战争结束后处理使节、边防和西藏军政善后。从地方治理、身份、家庭与物质生活的角度，这一节点要求追问它在何处改变了资源、道路、组织或日常风险。核心来源为《清高宗实录》乾隆五十七年相关条；《筹办夷务始末》相关年卷与中外照会、条约文本。它首先留下的是官府分类、任命、处置或资源调拨的痕迹；要判断基层是否照办，仍须寻找县档、账簿、契约或后续案件。账本所示的研究方向是条约文本、外交档案、口岸社会与国际法实践研究。可核的条件包括远征成本与边境稳定需要可持续安排。应分层观察的后果是朝贡表述与边境实际互动并不完全一致。证据界限：以乾隆五十七年（1792年）的同期文书为定位起点；官修记录、奏报或外交文本服务特定机构立场，具体日期、规模、地方执行与对手视角须以独立材料复核。
\hypertarget{social:EQ-1793-MACARTNEY-MISSION}{}\paragraph{乾隆五十八年（1793年）}马戛尔尼使团抵达清廷，围绕贸易、礼仪和外交关系交涉。从地方治理、身份、家庭与物质生活的角度，这一节点要求追问它在何处改变了资源、道路、组织或日常风险。核心来源为《清高宗实录》乾隆五十八年相关条；《筹办夷务始末》相关年卷与中外照会、条约文本。它首先留下的是官府分类、任命、处置或资源调拨的痕迹；要判断基层是否照办，仍须寻找县档、账簿、契约或后续案件。账本所示的研究方向是条约文本、外交档案、口岸社会与国际法实践研究。可核的条件包括英国对华贸易与全球帝国扩张需求增长。应分层观察的后果是使团记录具有强烈英方视角，礼仪争议不应简化为“闭关”。证据界限：以乾隆五十八年（1793年）的同期文书为定位起点；官修记录、奏报或外交文本服务特定机构立场，具体日期、规模、地方执行与对手视角须以独立材料复核。
\hypertarget{social:EQ-1794-WHITE-LOTUS-OUTBREAK}{}\paragraph{乾隆五十九年（1794年）}白莲教相关起事在川楚陕地区扩展，长期战争开始。从地方治理、身份、家庭与物质生活的角度，这一节点要求追问它在何处改变了资源、道路、组织或日常风险。核心来源为《清高宗实录》乾隆五十九年相关条；军机处档、战事方略及地方奏报。这类编年材料可以标定朝廷获知和叙述事件的时间，却往往压低妇女、儿童、佃户、流民和非文书精英的行动。账本所示的研究方向是战争动员、军需、战场暴力与地方社会研究。可核的条件包括地方生计压力、宗教网络和治理失灵相互作用。应分层观察的后果是“教匪”标签不能替代对不同参与群体的分析。证据界限：以乾隆五十九年（1794年）的同期文书为定位起点；官修记录、奏报或外交文本服务特定机构立场，具体日期、规模、地方执行与对手视角须以独立材料复核。
\hypertarget{social:EQ-1795-QIANLONG-ABDICATION}{}\paragraph{乾隆六十年（1795年）}乾隆帝禅位于嘉庆帝，自为太上皇，仍保留重要政治影响。从地方治理、身份、家庭与物质生活的角度，这一节点要求追问它在何处改变了资源、道路、组织或日常风险。核心来源为《清高宗实录》乾隆六十年相关条；宫中朱批奏折、内阁大库题本与《清史稿》相关志传。它首先留下的是官府分类、任命、处置或资源调拨的痕迹；要判断基层是否照办，仍须寻找县档、账簿、契约或后续案件。账本所示的研究方向是宫廷政治、制度运作与中央—地方信息链研究。可核的条件包括乾隆希望履行不逾祖父在位年数的承诺。应分层观察的后果是嘉庆初年权力交接受太上皇与和珅网络制约。证据界限：以乾隆六十年（1795年）的同期文书为定位起点；官修记录、奏报或外交文本服务特定机构立场，具体日期、规模、地方执行与对手视角须以独立材料复核。
